\begin{problem}{크림 먼저, 잼 먼저}{standard input}{standard output}{1 second}{256 megabytes}

영국에서 스콘을 맛있게 먹는 방법은 위에 크림을 바른 후 잼을 얹는 것, 혹은 위에 잼을 바른 후 크림을 얹는 것이다. 성준이는 파티에서 어느 방식의 스콘이 더 많은지 세어 보기로 했다. 불행하게도, 모든 스콘들은 하나의 탑으로 쌓여 있어, 위에서부터 재료를 읽을 수밖에 없었다.

스콘 탑의 재료들을 나타내는 문자열 $T$가 주어지면, 크림이 위인 스콘의 개수와, 잼이 위인 스콘의 개수를 각각 구하여라

문자열 $T$는 문자 $C$, $J$, $S$ 로만 이루어져 있다. $C$는 크림을, $J$는 잼을, $S$는 스콘을 나타낸다.

\begin{itemize}
  \item \textbf{크림이 위인 스콘}: 1개 이상의 $C$, 1개 이상의 $J$, 하나의 $S$가 연속해서 나타나는 경우이다.
  \item \textbf{잼이 위인 스콘}: 1개 이상의 $J$, 1개 이상의 $C$, 하나의 $S$가 연속해서 나타나는 경우이다.
  \item $C$와 $J$는 모두 자신의 오른쪽에서 가장 처음 오는 $S$에 바르는 것으로 가정한다. 즉, $CJCS$에서 $S$에는 $C$, $J$, $C$가 모두 발라져 있다.
\end{itemize}

만약 크림이나 잼 중 하나라도 빠져 있거나, 두 재료가 번갈아 여러 번 나타나는 스콘은 잘못된 스콘으로 간주하며, \textbf{어느 쪽에도 세지 않는다.}



\InputFile
첫째 줄에 문자열 $T$의 길이를 나타내는 정수 $N$이 주어진다.

둘째 줄에 문자열 $T$가 주어진다.

\OutputFile
첫째 줄에 크림이 위인 스콘의 개수를 출력한다.

둘째 줄에 잼이 위인 스콘의 개수를 출력한다.

\Examples

\begin{example}
\exmpfile{example.01}{example.01.a}%
\exmpfile{example.02}{example.02.a}%
\exmpfile{example.03}{example.03.a}%
\exmpfile{example.04}{example.04.a}%
\end{example}

\end{problem}

